\subsection{The cost of retaining the maximum witness count}

On auxiliary squares, the constructions above retain $b$ witness vertices.
For a longer auxiliary cycle, requiring the same witness count can lose an
unbounded number of vertices, even in connected templates of minimum degree
two and maximum degree three. The next family quantifies this loss exactly.

\begin{theorem}\label{thm:maximum-witness-loss}
For every integer $r\ge1$ there is a two-connected bipartite partial template
$H_r$ on $20r$ vertices, with $b=4r$ helpers, minimum degree two and maximum
degree three, such that
\[
 \beta_I(H_r)=11r+\lfloor r/2\rfloor.
\]
Among dissociation sets containing every helper and exactly $b$ witness
vertices, the largest size is $10r$. Thus this additional restriction loses
exactly $r+\lfloor r/2\rfloor$ vertices.
\end{theorem}
\begin{proof}
Use the index set $V=\{0,\ldots,8r-1\}$ and let
\[
 P=\{\{2j,2j+1\}:0\le j<4r\},\qquad
 Q=\{\{2j+1,2j+2\}:0\le j<4r\},
\]
where indices in $Q$ are reduced modulo $8r$. Thus $A=P\cup Q$ is the
cycle $C_{8r}$. Split $V$ into intervals $K_i=\{8i,\ldots,8i+7\}$,
$0\le i<r$. Within each $K_i$, put an owner four-cycle on the even offsets
$0,2,4,6$ and another on the odd offsets $1,3,5,7$, in those orders.
These are all edges of $B$.

Give index $v$ the physical witness colour $\lfloor v/2\rfloor\bmod2$.
Then $P$ joins equal colours, while $Q$ and $B$ join opposite colours.
Construct the physical template as before. Its helper--witness subgraph is
obtained by subdividing the $P$ edges of $C_{8r}$, so is connected. Each
owner has its witness neighbour, hence the whole graph is connected. The
helpers have degree two, and all witnesses and owners have degree three.
The specified colours, reversed on owners and helpers, give a bipartition.
The graph is two-connected. Deleting a helper or witness leaves the backbone
cycle as a connected path; every owner cycle still has at least three
witness attachments to that path. Deleting an owner leaves its owner cycle
as a path attached to the intact backbone. Thus every single vertex deletion
leaves a connected graph.

Fix any selected witness set $T$, which must be independent in $A$.
For $p\in\{0,1\}$ let $k_i^p$ count the selected witnesses of parity $p$
in $K_i$. In a four-cycle with $k$ vertices forbidden, the largest
dissociation set has size $\min\{2,4-k\}$: at most two vertices may be
selected, and every set of at most two vertices is dissociation. The owner
cycles are disjoint, so the largest physical set with this fixed $T$ has size
\[
 4r+\sum_{i,p}\bigl(k_i^p+\min\{2,4-k_i^p\}\bigr)
 =8r+\sum_i g_i,\qquad
 g_i=\min\{k_i^0,2\}+\min\{k_i^1,2\}.
\]
Here owners indexed by $T$ are forbidden, and all other owner choices are
independent between cycles. This maximization is therefore exact.

Before the cyclic joins are added, each $K_i$ is a path of order eight.
Consequently
$|T\cap K_i|\le4$ and $g_i\le4$. If $g_i=4$, then $T\cap K_i$ has exactly
two vertices of each parity. Both boundary offsets $0$ and $7$ must belong
to it: omitting $0$ leaves a path of order seven whose unique independent
four-set consists of the odd offsets, and omitting $7$ forces the even
offsets instead. Adjacent intervals cannot both have $g_i=4$, because the
last vertex of one is adjacent to the first of the next. For $r=1$ the
closing edge similarly forbids $g_0=4$; for $r\ge2$, at most
$\lfloor r/2\rfloor$ intervals can have that value. Every other interval
contributes at most three. Hence
\[
 \beta_I(H_r)\le8r+3r+\lfloor r/2\rfloor.
\]

For equality choose $\lfloor r/2\rfloor$ pairwise nonadjacent intervals in
the cyclic order, choosing none when $r=1$. In each chosen interval select
witness offsets $\{0,2,5,7\}$ and owner offsets $\{1,3,4,6\}$; in each
other interval select witness offsets $\{1,3,6\}$ and owner offsets
$\{0,2,5,7\}$. The witnesses are independent, including across interval
boundaries. Each owner cycle receives exactly two vertices and no selected
owner shares its index with a selected witness. Together with all helpers
these choices form a dissociation set of the asserted size.

Finally, an independent set of maximum size $4r$ in $C_{8r}$ is one of its
two bipartition classes. Either witness choice forbids all four vertices
of one owner cycle in each interval and none of the other cycle. Exactly
$2r$ owners can then be selected. Including the $4r$ helpers gives the
restricted optimum $4r+4r+2r=10r$, attained by choosing two owners in each
available four-cycle. Subtracting gives the stated loss.
\end{proof}

This family concerns the restriction on the witness count. It satisfies the
earlier $8b/3$ bound and has no auxiliary four-cycle component, so
Theorem~\ref{thm:owner-matching} applies. No unrestricted dissociation optimum
or new subcubic gap coefficient is asserted. To combine square and longer
auxiliary components, a construction must allow the witness count to change;
even a single long auxiliary cycle can require such a change.
